\documentclass{resume}
\usepackage{zh_CN-Adobefonts_external} 
\usepackage{linespacing_fix}
\usepackage{cite}
\usepackage{hyperref}
\hypersetup{
    colorlinks=true,
    linkcolor=cyan,
    filecolor=magenta,      
    urlcolor=blue,
}

\begin{document}
\pagenumbering{gobble}

%***"%"后面的所有内容是注释而非代码，不会输出到最后的PDF中
%***使用本模板，只需要参照输出的PDF，在本文档的相应位置做简单替换即可
%***修改之后，输出更新后的PDF，只需要点击Overleaf中的“Recompile”按钮即可

%在大括号内填写其他信息，最多填写4个，但是如果选择不填信息，
%那么大括号必须空着不写，而不能删除大括号。
%\otherInfo后面的四个大括号里的所有信息都会在一行输出
%如果想要写两行，那就用两次这个指令（\otherInfo{}{}{}{}）即可


%***********个人信息**************
\MyName{XXX}
\sepspace
\SimpleEntry{\textbf{籍贯：}}
\SimpleEntry{\textbf{邮箱：}}
\SimpleEntry{\textbf{电话：}}
\SimpleEntry{\textbf{政治面貌：}}
\SimpleEntry{\textbf{通讯地址：}}
\SimpleEntry{\textbf{个人主页：}\href{https://doublelll3.ml}{https://doublelll3.ml}}

%************照片**************
%照片需要放到images文件夹下，名字必须是you.jpg，注意.jpg后缀（可以去resume.cls第101行处修改），如果不需要照片可以不添加此行命令
%0.15的意思是，照片的宽度是页面宽度的0.15倍，调整大小，避免遮挡文字
\yourphoto{0.13}

%***********教育背景**************
\section{教育背景}
%***第一个大括号里的内容向左对齐，第二个大括号里的内容向右对齐
%***\textbf{}括号里的字是粗体，\textit{}括号里的字是斜体
\datedsubsection{\textbf{xxxx大学}，xxxxx及控制，\textit{硕士}}{2018.9 - 2021.7}
\datedsubsection{\textbf{xxxx大学}，xxxxxx及控制，\textit{本科}}{2014.9 - 2018.7}


%***********实习经历**************
\section{实习经历}
\datedsubsection{\textbf{xxxx公司（实习）}，xx工程师}{2020.6 - 2020.9}
\Content
{针对xxxx，实时检测xxxx情况（主要负责人）。}
{研究xxx算法、xxxx算法以及两者的结合。}
{完成xxxx与xxx的结合，在xxxx，目前已上线使用。}

\datedsubsection{\textbf{xxx研究院（实习）}，xxxx实习生}{2020.1 - 2020.5}
\Content
{负责xxx项目，对xxx进行智能诊断。}
{搭建xxx网站，结合xxx，利用xxx，基于xxx实现xxx。}
{1.完成xxx；2.基于xxx，算法准确率在85\%左右。}

\datedsubsection{\textbf{xxx大学研究生课题}，xxx}{2018.3 - 至今}
\Content
{针对xxx的不足，研究一种基于xxx对xxx进行xxx的方法。}
{基于xxx建立xxx模型，使xxx，基于xxx，设计一种xxx。}
{1.建立了xxx，并在xxx得到了验证；2.对xxx，算法准确率在95\%以上；3.对xxx可以自主学习，从而xxx。}

\datedsubsection{\textbf{大学生创新创业项目}，xxx}{2016.3 - 2017.3}
\Content
{针对xxx，基于xxx，加入xx实现了xxx的功能（主要负责人）。}
{搭建xxx，利用xxx的输出，实现对xxx的控制。}
{1.实现了xxx；2.获得了xxx奖等荣誉。}

%***********项目经历**************
\section{项目经历}
\datedsubsection{\textbf{xxxx（项目）}，xxxx}{2020.6 - 2020.9}
\Content
{针对xxxx，实时检测xxxx情况（主要负责人）。}
{研究xxx算法、xxxx算法以及两者的结合。}
{完成xxxx与xxx的结合，在xxxx，目前已上线使用。}


%***********学生工作**************
\section{学生工作}

\datedsubsection{\textbf{校学生会}，宣传部长，\textit{负责学生会整体宣传策划与执行，统筹多渠道传播运营}}{2023.9 - 2025.6}

\datedsubsection{\textbf{校足球队}，队长，\textit{负责队伍日常管理与训练组织，统筹赛事筹备与场上指挥}}{2022.9 - 2026.6}

\datedsubsection{\textbf{创新协会}，组织部长，\textit{负责协会活动组织策划与落地执行，统筹跨部门协作与成员管理}}{2023.9 - 2025.6}

%***********个人评价**************
\section{个人评价}

{CET-6 542分，雅思7.0，能独立进行英文技术交流与文档写作，每周精读英文技术资料3-5篇。}

{主力Python（3年，30k+行），掌握Java、TypeScript，代码规范严格，单元测试覆盖率维持65\%-80\%。}

{xxx。}

{xxx}

{xxx}

{xxx}

{xxx}

{每周跑步10-15km；维护技术知识库笔记420+条，每月产出技术总结，持续学习成长。}

\end{document}

